% ============================================================ %
% Runtime Monitoring over Signed-Incidence Hypergraphs           %
% Paper skeleton targeting RV 2026 (Springer LNCS, 12 pages)     %
% Backup venues: FormaliSE, FMICS, FM tool papers track          %
% ============================================================ %

\documentclass[runningheads]{llncs}
\usepackage{amsmath,amssymb,amsthm}
\usepackage{graphicx}
\usepackage{listings}
\usepackage{xcolor}
\usepackage{hyperref}

\begin{document}

\title{Streaming Runtime Monitoring of Signal Temporal Logic
       over Signed-Incidence Hypergraphs}
\titlerunning{Runtime Monitoring over Signed Hypergraphs}
\author{Csaba Hajdu\inst{1,2}}
\institute{Sz\'echenyi Istv\'an University, Gy\H{o}r, Hungary \and
           SZTIP, Sz\'ekesfeh\'erv\'ar, Hungary \\
           \email{hajdu.csaba@ga.sze.hu}}
\maketitle

% ---------- Abstract ---------- %

\begin{abstract}
% ~150 words, last.
% STRUCTURE to fill in at the end:
%   (1) context: runtime verification of CPS requires structural
%       state, not only scalar signals.
%   (2) gap: STL/LTL/CTL are defined over scalar traces or Kripke
%       structures; directed hypergraphs with signed incidence are
%       the natural substrate for multi-entity, multi-context
%       constraints but lack temporal-logic semantics.
%   (3) contribution: (a) formal semantics of STL over signed-
%       incidence hypergraph traces, with structural robustness;
%       (b) bounded-memory online monitor construction; (c) Rust
%       implementation on top of the HyMeKo IR; (d) case study
%       on cross-layer kinematic/safety constraints.
%   (4) result: monitor runs in X microseconds per observation on
%       Y-vertex hypergraphs, detects Z structural violations that
%       classical scalar STL monitors cannot express.
\keywords{Runtime verification \and Signal temporal logic \and
          Hypergraphs \and Cyber-physical systems \and HyMeKo.}
\end{abstract}

% ============================================================ %
\section{Introduction}
\label{sec:intro}
% Target: 1 page. Three paragraphs + contribution list.
% ============================================================ %

% Paragraph 1: The runtime-verification gap for CPS.
% - Classical RV: Maler-Nickovic online STL, Havelund-Rosu
%   rewriting, Bauer-Leucker-Schallhart impartiality.
% - All operate on scalar signal traces.
% - Modern CPS (industrial robotics, autonomous vehicles,
%   multi-agent systems) have constraints that are structural,
%   not scalar: "every joint has a corresponding actuator in
%   the runtime stack", "no kinematic cycle is present in the
%   live model", "maintenance-derived health state constrains
%   permissible production speed". These properties are over
%   the running system's structure and its cross-layer
%   dependencies, not over isolated signals.
% - Cite: Maler-Nickovic 2004, Donze-Maler 2010, Bauer et al.
%   2011, Bartocci et al. RV survey 2018.

% Paragraph 2: Why hypergraphs.
% - Signed-incidence directed hypergraphs are the substrate
%   used by HyMeKo (cite) for cross-view consistency in CPS
%   engineering. n-ary relations stay n-ary, signs disambiguate
%   directional roles, and the Levi bipartite presentation
%   (Levi 1942, Berge 1973, Gallo et al. 1993) gives a uniform
%   traversal substrate.
% - But: HyMeKo compiles at design time. A running CPS needs
%   monitoring at runtime: structural evolution observed,
%   temporal-logic property evaluated incrementally, verdict
%   produced within bounded latency.

% Paragraph 3: This paper.
% - We lift LTL, CTL, and STL onto signed-incidence hypergraph
%   traces. Atomic propositions are predicates over the
%   hypergraph (using HyMeKo's existing predicate algebra).
%   Temporal operators apply over sequences of hypergraph
%   states.
% - We give a bounded-memory online monitor construction for
%   the bounded-horizon fragment of STL. The monitor runs
%   alongside HyMeKo and observes attribute and incidence
%   updates, producing robustness verdicts incrementally.
% - We demonstrate on two case studies: kinematic-limit
%   monitoring and cross-layer constraint propagation (the
%   Szilagyi et al. scenario, cite).

\paragraph{Contributions.}
\begin{enumerate}
  \item A formal semantics for STL (and, by restriction, LTL
        and qualitative CTL) over traces of signed-incidence
        directed hypergraphs, including a \emph{structural
        robustness} function that extends the Maler-Nickovic
        robustness to predicates over hypergraph structure.
  \item A bounded-memory online monitor construction for the
        bounded-horizon fragment of hypergraph STL, with
        complexity $\mathcal{O}(|\phi| \cdot k)$ per observation
        where $k$ is the sliding-window length determined by
        the formula's temporal horizon, leveraging incremental
        predicate evaluation (cf.\ Rath et al.\ 2008).
  \item An open-source Rust implementation (\texttt{hymeko\_monitor})
        built on top of the HyMeKo IR.
  \item Empirical validation on a kinematic-chain limit
        scenario and a cross-layer health-to-speed constraint
        scenario drawn from industrial robotic ecosystems.
\end{enumerate}

% ============================================================ %
\section{Background}
\label{sec:background}
% Target: 1.25 pages. Four subsections, each ~0.3 page.
% ============================================================ %

\subsection{Signed-Incidence Directed Hypergraphs}
% H = (V, E, I, \tau, \sigma, \alpha). Levi bipartite presentation.
% Cite HyMeKo, CogInfoCom 2024, Levi 1942, Berge 1973, Gallo et al. 1993.
% Concrete example: 3-link chain with 2 joint hyperedges and
% signed parent/child incidences. One figure, half-column.

\subsection{Signal Temporal Logic}
% Classical syntax: \phi ::= p | \neg\phi | \phi \wedge \psi |
% F_{[a,b]}\phi | G_{[a,b]}\phi | \phi U_{[a,b]} \psi.
% Boolean and robust semantics per Maler-Nickovic 2004, Donze-Maler 2010.
% Robustness: \rho(p, s, t) = \mu(s(t)), where \mu is a signed
% margin from the predicate threshold.
% Temporal combinators: sup/inf over the interval.

\subsection{Runtime Verification}
% Online monitors, impartiality, 3-valued LTL verdicts
% (Bauer-Leucker-Schallhart 2011), STL online monitoring with
% predictive windows (Maler-Nickovic, Deshmukh et al. 2017).
% VIATRA-CEP (Daboczi et al. 2017) as the closest prior work on
% graph-structured runtime verification -- but typed graphs,
% not hypergraphs, and boolean-only.

\subsection{Incremental Pattern Matching}
% Rath-Bergmann-Okros-Varro 2008 and follow-up. Rete-style
% dependency networks, delta-propagation of updates through
% partial matches. We extend this to hypergraph predicates
% with signed incidences in Section~\ref{sec:monitor}.

% ============================================================ %
\section{Semantics of Temporal Logic over Signed Hypergraphs}
\label{sec:semantics}
% Target: 3 pages. The technical core of the paper. This section
% must be tight, with numbered definitions and propositions.
% ============================================================ %

\subsection{Hypergraph Traces}
% DEFINITION 1 (Hypergraph state). Same as the HyMeKo IR tuple.
% DEFINITION 2 (Trace). A trace is a finite or infinite sequence
%   \pi = (H_0, t_0), (H_1, t_1), ... with t_i \in \mathbb{R}_{\geq 0},
%   strictly increasing. Piecewise-constant interpretation:
%   H(t) = H_i for t \in [t_i, t_{i+1}).
% DEFINITION 3 (Transition kinds). Two families:
%   (a) attribute update: \alpha changes, rest fixed.
%   (b) structural rewrite: V, E, I, \tau, \sigma may change.
%   This paper focuses on (a); (b) is future work.

\subsection{Structural Atomic Propositions}
% DEFINITION 4 (Hypergraph predicate). A predicate is a function
%   P: H -> \mathbb{B} defined by a closed expression over the
%   HyMeKo predicate algebra extended with:
%     COUNT(p) -- number of entities satisfying p
%     EXISTS(p), FORALL(p)
%     ARITY(e) -- |{v : (v,e) \in I}|
%     ATTR(v, k) -- \alpha(v)[k] read as scalar
%     ACYCLIC(), CONNECTED(), ...
% DEFINITION 5 (Robustness of a predicate). For predicates with
%   natural numeric content, robustness is the signed margin:
%     \rho(COUNT(p) \geq n, H) = |{v : P(v)}| - n
%     \rho(ATTR(v, k) \leq \theta, H) = \theta - \alpha(v)[k]
%     \rho(ARITY(e) = d, H) = -(|incid(e)| - d)^2     [or abs]
%   For structural-only predicates (ACYCLIC, CONNECTED),
%   robustness is \pm\infty (qualitative case).
% PROPOSITION 1. Robustness is sound with respect to boolean
%   evaluation: \rho(P, H) > 0 \iff P(H) = true, and \rho(P, H)
%   < 0 \iff P(H) = false, with \rho = 0 on the boundary.

\subsection{STL over Hypergraph Traces}
% SYNTAX (as in classical STL).
% BOOLEAN SEMANTICS: standard, over \pi, t.
% ROBUST SEMANTICS:
%   \rho(P, \pi, t)      = \rho(P, H(t))
%   \rho(\neg \phi)      = -\rho(\phi)
%   \rho(\phi \wedge \psi) = min(\rho(\phi), \rho(\psi))
%   \rho(F_{[a,b]} \phi, \pi, t) = sup_{t' \in [t+a, t+b]} \rho(\phi, \pi, t')
%   \rho(G_{[a,b]} \phi, \pi, t) = inf_{t' \in [t+a, t+b]} \rho(\phi, \pi, t')
%   \rho(\phi U_{[a,b]} \psi, \pi, t) = sup_{t' \in [t+a, t+b]}
%                                        min(\rho(\psi, \pi, t'),
%                                            inf_{t'' \in [t, t']} \rho(\phi, \pi, t''))
% PROPOSITION 2 (soundness). \rho(\phi, \pi, t) > 0 \iff
%   (\pi, t) \models \phi.

\subsection{LTL and CTL as Fragments}
% LTL: STL with intervals [0, \infty) and qualitative predicates.
% CTL: requires computation tree, not trace; appears in future
%   work when structural rewrites (transition kind b) are
%   included. Section ends with: "for the remainder of the
%   paper, we focus on STL over attribute-update traces."

% ============================================================ %
\section{Monitor Construction}
\label{sec:monitor}
% Target: 2 pages. Algorithm + complexity.
% ============================================================ %

\subsection{Bounded-Horizon Fragment}
% DEFINITION 6. A formula is bounded-horizon if every temporal
%   operator has a bounded interval [a, b] with b < \infty.
%   This fragment admits bounded-memory online monitoring.

\subsection{Sliding-Window Monitor}
% ALGORITHM:
%   (1) Parse \phi into a tree.
%   (2) Compute per-subformula horizon H_\phi = max b over
%       temporal operators above the subformula.
%   (3) Allocate a ring buffer of size ceil(H_\phi / dt_min)
%       per subformula.
%   (4) On each sample (H_i, t_i):
%       (a) For each predicate p in \phi:
%            -- if p's match set could have changed given the
%               incremental delta from H_{i-1} to H_i, re-evaluate
%               robustness and push into p's buffer.
%       (b) For each temporal subformula, recompute its
%           robustness via sup/inf over its window.
%       (c) Emit \rho(\phi, \pi, t_i - H_\phi) as the
%           \emph{settled} verdict at the window trailing edge.

\subsection{Incremental Predicate Evaluation}
% Extend Rath et al.'s dependency tracking: each predicate
% caches the set of vertices / hyperedges it matched on in the
% previous state. An attribute update on vertex v triggers
% re-evaluation only of predicates whose match set includes v.
% Structural predicates (ACYCLIC, CONNECTED) fall back to
% global re-evaluation; this is the cost paid for their
% expressive power.

\subsection{Complexity}
% PROPOSITION 3. For bounded-horizon \phi with |\phi|
%   subformulas, horizon H_\phi, minimum sample spacing
%   dt_min, and hypergraph size |V| + |E|:
%   per-observation cost is
%     O(|\phi| \cdot H_\phi / dt_min + \Delta)
%   where \Delta is the incremental delta between successive
%   hypergraph states. For attribute-only updates with a
%   small delta, \Delta is dominated by affected predicate
%   match sets.

% ============================================================ %
\section{Implementation}
\label{sec:impl}
% Target: 1 page. Crate architecture + API example.
% ============================================================ %

% - \texttt{hymeko\_monitor} Rust crate (~2000 LOC).
% - Crates: hymeko\_core (IR), hymeko\_monitor (this work).
% - Architecture diagram (Fig.\ \ref{fig:arch}).
% - Public API listing (one short Rust listing, ~15 lines,
%   showing formula construction + monitor stepping).
% - Released as open source; code, fixtures, and benchmarks in
%   the companion repository (cite).

% ============================================================ %
\section{Case Studies}
\label{sec:cases}
% Target: 2 pages. Two scenarios, empirical numbers.
% ============================================================ %

\subsection{Case Study 1: Kinematic Limit Monitoring}
% SETUP: anthropomorphic-arm fixture (7 links, 7 joints), with
% joint positions attributed to joint hyperedges. Property:
%   G(collaborative_mode -> G(FORALL joint. WITHIN_LIMITS(joint))).
% TRACE: 30 s simulated motion, sampling at 100 Hz (3000 states).
% RESULT: monitor detects first limit approach at t=X with
% robustness crossing zero; violation at t=Y with negative
% robustness. Wall-clock overhead < Z microseconds per sample
% on the fixture.

\subsection{Case Study 2: Cross-Layer Health-to-Speed Constraint}
% SETUP: reproduce the Szilagyi et al. maintenance/production/
% safety scenario (cite). Property:
%   G((HEALTH < theta_h) -> G_{[0, 100ms]}(SPEED <= v_maint_limit)).
% The property spans two context subgraphs (maintenance D_out
% and production D_out) and requires the monitor to detect
% violations of cross-layer constraint propagation within a
% bounded latency.
% RESULT: monitor detects the delayed-response failure mode;
% reports robustness margin live during correct operation.

\subsection{What Classical STL Cannot Express}
% Short paragraph: the closest scalar-STL formulation requires
% pre-aggregating hypergraph structure into scalar features,
% which loses the cross-layer dependency structure. Demonstrate
% with a specific property that is expressible in hypergraph
% STL and is not expressible (or requires multiple coupled
% scalar monitors) in classical STL.

% ============================================================ %
\section{Related Work}
\label{sec:related}
% Target: 0.75 page.
% ============================================================ %

% - Classical RV: Bauer-Leucker-Schallhart 2011; Bartocci et al.
%   RV survey 2018; Falcone et al. RV tutorial.
% - STL: Maler-Nickovic 2004; Donze-Maler 2010 (robustness);
%   Deshmukh et al. 2017 (predictive monitoring).
% - Graph-structured RV: VIATRA-CEP (Daboczi et al. 2017);
%   Rath et al. 2008 (incremental pattern matching, the
%   algorithmic ancestor here).
% - STL-as-reward in RL: Aksaray et al. 2016 (Q-learning with
%   STL), Balakrishnan-Deshmukh 2019, the Belta group line,
%   Hasanbeig et al. (LTL rewards). We connect to this line in
%   Section~\ref{sec:conclusion} as future work.
% - Hypergraph learning: Feng et al. HGNN 2019, Bai et al. 2021.
%   None address runtime verification.

% ============================================================ %
\section{Limitations and Future Work}
\label{sec:future}
% Target: 0.5 page.
% ============================================================ %

% - Transition kind (b): structural rewrites introduce a
%   computation tree and enable CTL monitoring. Nontrivial
%   because the rewrite relation is HyMeKo's rewriting engine;
%   the monitoring question becomes model-checking-like.
% - Robustness design for purely structural predicates: current
%   \pm\infty choice is sound but does not give dense reward
%   gradients. Future work: quantitative distance metrics on
%   hypergraph structure (edit distance from violating pattern).
% - RL integration: robustness as dense reward signal, shield
%   synthesis from CTL invariants. Part of the broader program
%   sketched in [cite: Szilagyi et al.].
% - Distributed monitoring: split the monitor across context
%   subgraphs for decentralized CPS.

% ============================================================ %
\section{Conclusion}
\label{sec:conclusion}
% Target: 0.25 page. One paragraph.
% ============================================================ %

% Recap contributions in plain language. Emphasize: this is the
% first temporal-logic runtime monitor for directed hypergraphs
% with signed incidence. The semantics, the monitor construction,
% and the Rust implementation together provide a deployable
% substrate for structural runtime verification in CPS. The RL
% reward-shaping connection is the next step.

% ============================================================ %
\begin{thebibliography}{99}
% Required citations, to populate:
% - Maler \& Nickovic 2004 (STL)
% - Donze \& Maler 2010 (robustness)
% - Bauer-Leucker-Schallhart 2011 (3-valued LTL)
% - Bartocci et al. 2018 (RV survey)
% - Deshmukh et al. 2017 (online STL)
% - Rath-Bergmann-Okros-Varro 2008 (incremental pattern matching)
% - Daboczi et al. 2017 (VIATRA-CEP)
% - Aksaray et al. 2016 (Q-learning with STL)
% - Balakrishnan \& Deshmukh 2019 (STL-RL)
% - Levi 1942, Berge 1973, Gallo et al. 1993 (hypergraphs)
% - Hajdu \& Csapo CogInfoCom 2024 (HyMeKo lineage)
% - Hajdu et al. SMC 2026 (HyMeKo cross-view consistency)
% - Szilagyi et al. Technologies 2026 (the cross-layer scenario)
\end{thebibliography}

\end{document}
